%%
%% ACM MM sigconf format
%%
\documentclass[sigconf, anonymous, review]{acmart}
%%
%% Rights management information (placeholder - replace when you get the real one)
%%
% \setcopyright{acmlicensed}
% \copyrightyear{2025}
% \acmYear{2025}
% \acmDOI{XXXXXXX.XXXXXXX}
\acmConference[MM '26]{Proceedings of the 34th ACM International Conference on Multimedia}{November 10--14, 2026}{Rio de Janeiro, Brazil}
% \acmISBN{978-1-4503-XXXX-X/2025/10}
\setcopyright{none}
\settopmatter{printacmref=false}
\renewcommand\footnotetextcopyrightpermission[1]{}
%%
%% Additional packages (compatible with acmart)
%%
\usepackage{algorithm}
\usepackage{algorithmic}
\usepackage{amsmath}
\usepackage{amssymb}
\usepackage{amsthm}
\usepackage{booktabs}
\usepackage{enumitem}
\usepackage{graphicx}
\usepackage{multirow}
\usepackage{subfig}
\usepackage{array}
\newcolumntype{C}{>{\centering\arraybackslash}m{4em}}

\acmSubmissionID{6807}

\begin{document}

%%
%% Title
%%
\title{CAGRec: Collaborative-Anchor Guided Multimodal Generative Recommender}

%%
%% Authors - TODO: replace with your own info
%%

%%
%% Abstract
%%
\begin{abstract}
Multimodal Generative Recommendation (MGR) encodes items as discrete semantic identifiers (IDs) via vector quantization and predicts the next item through autoregressive token generation. 
% 
Conventional multi-modal approaches construct item codes from text and image features through cross-modal RQ-VAE, consistently outperforming single-modality methods. 
% 
However, existing methods build code spaces purely from content features, overlooking that frequently co-occurring items in user interaction sequences can differ greatly in content.
% 
Moreover, autoregressive decoding ranks candidates solely by token-level likelihood, unable to evaluate them in continuous semantic space. 
% 
We propose Collaborative-Anchor guided multimodal Generative Recommender (CAGRec) to address both challenges. 
% 
For the content-collaborative gap, we anchor the cross-modal RQ-VAE with collaborative embeddings derived from user interaction sequences, steering both text and image code spaces toward behavioral structure.
%
% The anchored RQ-VAE decoder gives rise to a collaborative-semantic space that jointly encodes behavioral and content signals. Existing contrastive alignment objectives, when operating in this enriched space, capture fine-grained collaborative similarity that code-level supervision alone cannot express.
% [已确认: 老师批准修改摘要，以下为新版本，显式提及 CAR]
We further propose Collaborative-Aware Regularization~(CAR), a contrastive loss that clusters collaboratively related items in the latent space before code assignment, ensuring behavioral proximity survives quantization.
%
The anchored RQ-VAE decoder gives rise to a collaborative-semantic space that jointly encodes behavioral and content signals.
%
For the discrete-continuous gap, we design a semantic-aware reranking strategy that combines generation likelihood with collaborative-semantic similarity, pulling candidate evaluation back into the continuous space.
% 
Extensive experiments on three recommendation datasets show that CAGRec consistently outperforms state-of-the-art baselines.
% 
Our codes are available at \url{https://anonymous.4open.science/r/CAGRec}.
\end{abstract}

%%
%% CCS Concepts - TODO: generate your own at https://dl.acm.org/ccs/ccs.cfm
%%
\begin{CCSXML}
<ccs2012>
 <concept>
  <concept_id>10002951.10003317.10003347.10003350</concept_id>
  <concept_desc>Information systems~Recommender systems</concept_desc>
  <concept_significance>500</concept_significance>
 </concept>
 <concept>
  <concept_id>10010147.10010178.10010179</concept_id>
  <concept_desc>Computing methodologies~Natural language processing</concept_desc>
  <concept_significance>300</concept_significance>
 </concept>
</ccs2012>
\end{CCSXML}

\ccsdesc[500]{Information systems~Recommender systems}
% \ccsdesc[300]{Computing methodologies~Natural language processing}

%%
%% Keywords
%%
\keywords{Generative Recommendation, Multimodal Learning, Cross-modal Quantization, Semantic Identifiers}

\maketitle

%%%%%%%%%%%%%%%%%%%%%%%%%%%%%%%%%%%%%%%%%%%%%%%%%%%%%%%%%%%%%%%%%%%%%%%%
%% INTRODUCTION
%%%%%%%%%%%%%%%%%%%%%%%%%%%%%%%%%%%%%%%%%%%%%%%%%%%%%%%%%%%%%%%%%%%%%%%%

\begin{figure}[t]
    \centering
    \includegraphics[width=\linewidth]{Intro.pdf}
    \caption{Illustrative comparison on the Instruments dataset. (a)~Collaborative embeddings are extracted from the user interaction history to anchor item representations. (b)~Content-only methods rank the ground truth (amplifier) last. (c)~With collaborative anchoring, the amplifier rises to rank~\#1 and the model surfaces items (effects pedal) that content-only methods miss.}
    \Description{Illustrative comparison: content-only methods rank the ground truth item last; after anchoring item representations with collaborative embeddings extracted from the interaction history, the ground truth rises to rank one.}
    \label{fig:intro}
\end{figure}
\section{Introduction}
Recommendation systems help users navigate information overload across e-commerce~\cite{amazonrec1,taobaorec2}, social media~\cite{davidson2010youtube,dnn}, and news services~\cite{news_rec}.
%
With the rapid progress of large language models (LLMs), generative recommendation (GR) has emerged as a promising paradigm~\cite{wang2023generative,li2024large} that recasts recommendation as next-token prediction.
%
Each item is encoded as a discrete semantic identifier (ID) through vector quantization~\cite{vqvae}, and an autoregressive model generates the next item's semantic ID given the user's interaction history~\cite{tiger}.

A central question in GR is how to construct high-quality semantic IDs.
%
Early methods quantize single-modality (text) embeddings through RQ-VAE~\cite{image_rqvae,tiger}.
%
Later work improves tokenization via collaborative regularization~\cite{letter} or language model adaptation~\cite{LCRec}.
%
However, a single modality offers limited discriminability.
%
Recent multi-modal methods integrate text and visual features to produce richer item codes.
%
MMGRec~\cite{liu2024mmgrec} uses Graph RQ-VAE to generate item codes from multi-modal features propagated through user-item interaction graphs.
%
MQL4GRec~\cite{mql} builds a shared quantitative language across modalities and domains.
%
MACRec~\cite{macrec} applies cross-modal contrastive learning during residual quantization to reduce codebook collisions and improve code quality.
%
These multi-modal extensions consistently outperform single-modality methods.

Despite this progress, two fundamental challenges remain unresolved:

\noindent\textbf{Challenge 1: Content-Collaborative Gap.}
Existing multi-modal GR methods build code spaces primarily from content features without explicitly aligning the code structure with collaborative structure.
%
Content similarity, however, does not equal collaborative similarity.
%
As shown in Figure~\ref{fig:intro}, given a user who purchased an electric guitar, a tuner, and a guitar pick, content-only methods assign the highest scores to a guitar cable (0.82) and a guitar case (0.79)---the cable's description shares instrument-accessory keywords with the guitar, and the case's guitar-shaped silhouette yields high visual similarity---while the ground truth amplifier receives only 0.41 and ranks last, because its description is dominated by audio-electronics specifications (wattage, channels, EQ, reverb) and its boxy form bears no visual resemblance to the guitar.
%
By extracting collaborative embeddings from the interaction history and anchoring item representations with these embeddings, the amplifier rises to rank~\#1 (score 0.91): users who buy guitars frequently purchase amplifiers next.
%
The model also surfaces an effects pedal (0.80) that content-only methods entirely miss.
%
When the code space primarily reflects content proximity, collaboratively related items receive distant codes, and the autoregressive decoder must predict across a structure that conflicts with the recommendation objective.

\noindent\textbf{Challenge 2: Discrete-Continuous Gap.}
At inference, autoregressive decoding ranks candidates solely by token-level log-probabilities in the discrete code space; two items with nearly identical behavioral patterns can receive very different beam scores if their codes diverge at early positions.

Several recent studies have tried to address this gap.
%
LETTER~\cite{letter} adds a collaborative regularization loss during item tokenization, but treats collaborative signals as an auxiliary objective without structurally reshaping the code structure, and is limited to a single content modality.
%
ColaRec~\cite{colarec} derives item identifiers from pretrained collaborative filtering models and aligns them with content representations through contrastive learning, yet also operates on a single modality and cannot exploit multi-modal complementarity.
%
UNGER~\cite{unger} integrates semantic and collaborative knowledge into a unified code through adaptive embedding learning, but compresses all information into one code, sacrificing per-modality discriminability.
%
None of these methods simultaneously anchors multi-modal code learning with collaborative structure while preserving per-modality code spaces, nor do they leverage collaborative-semantic similarity to refine candidate ranking at inference.
%
We propose \textbf{C}ollaborative-\textbf{A}nchor \textbf{G}uided Multimodal Generative \textbf{Rec}ommender~(\textbf{CAGRec}), which bridges multi-modal code learning with collaborative signals.
%
We anchor the cross-modal RQ-VAE with collaborative embeddings learned from user interaction sequences, injecting behavioral structure into both text and image code spaces so that collaboratively related items receive proximate codes.
%
To reinforce this structure in the learned latent space, we propose Collaborative-Aware Regularization~(CAR), a contrastive objective that clusters collaboratively related items before code assignment, ensuring behavioral proximity survives the quantization bottleneck.
%
The anchored RQ-VAE decoder gives rise to a collaborative-semantic space that jointly encodes behavioral and content signals; during training, we further propose Collaborative Preference Anchoring~(CPA), a cosine alignment objective that grounds user representations in this space.
%
At inference, we design a semantic-aware reranking strategy that combines beam search likelihood with collaborative-semantic similarity, pulling candidate evaluation back into the continuous space.

Our main contributions are threefold:
\begin{itemize}
    \item We propose collaborative-anchored cross-modal quantization, which concatenates collaborative embeddings with text and image features as the RQ-VAE input, together with Collaborative-Aware Regularization~(CAR) that reinforces collaborative clustering in the latent space before code assignment, jointly restructuring both code spaces to respect behavioral proximity while preserving per-modality discriminability.
    \item We design a semantic-aware reranking strategy that re-evaluates beam search candidates by combining generation likelihood with continuous collaborative-semantic similarity, bridging the discrete-continuous gap at inference.
    \item We conduct extensive experiments on three recommendation datasets, showing that CAGRec consistently outperforms state-of-the-art baselines.
\end{itemize}

%%%%%%%%%%%%%%%%%%%%%%%%%%%%%%%%%%%%%%%%%%%%%%%%%%%%%%%%%%%%%%%%%%%%%%%%
%% RELATED WORK
%%%%%%%%%%%%%%%%%%%%%%%%%%%%%%%%%%%%%%%%%%%%%%%%%%%%%%%%%%%%%%%%%%%%%%%%

\section{Related Work}

\subsection{Generative Recommendation}

Classical sequential recommenders---recurrent models \cite{gru4rec}, self-atten\-tive Transformers \cite{sasrec}, masked Transformers \cite{sun2019bert4rec}, and multi-modal variants \cite{wei2019mmgcn,wang2023missrec}---rank items by representation similarity in an embed-and-retrieve paradigm.
%
Generative recommendation (GR) departs from this paradigm: each item is encoded as a discrete semantic ID via RQ-VAE~\cite{vqvae,image_rqvae}, and a sequence-to-sequence model autoregressively generates the target ID from the user's interaction history.
%
TIGER~\cite{tiger} establishes this paradigm by quantizing single-modality content embeddings into hierarchical codes and decoding the next item's code with a Transformer.
%
LC-Rec~\cite{LCRec} further adapts LLMs for GR through alignment tuning that bridges the gap between language semantics and recommendation semantics.
%
These single-modality methods rely solely on text to construct codes; text captures brand or category information but cannot represent visual attributes such as color, shape, or style, limiting the discriminability of the resulting semantic IDs.
%
Multi-modal extensions address this by incorporating complementary content signals so that items indistinguishable in one modality can be separated in another.
%
MMGRec~\cite{liu2024mmgrec} propagates multi-modal features through user-item interaction graphs before quantization via Graph RQ-VAE.
%
MQL4GRec~\cite{mql} encodes text and image features into a shared quantitative language and transfers recommendation knowledge across domains through pretraining.
%
MACRec~\cite{macrec} applies cross-modal contrastive learning during residual quantization, reducing codebook collisions by exploiting inter-modality complementarity while maintaining separate text and image code paths.
%
These methods improve semantic ID quality through richer content signals, yet the code spaces themselves remain content-driven without explicit collaborative guidance.

\subsection{Collaborative and Content Signal Alignment}

Collaborative filtering methods~\cite{rendle2009bpr,koren2009matrix,he2017neural} have long shown that items with distinct descriptions and appearances can share similar interaction patterns; content similarity does not equal collaborative similarity.
%
Bridging this gap is critical for GR, because the code structure directly constrains autoregressive prediction---collaboratively related items that receive distant codes force the decoder to predict across a conflicting structure.

LETTER~\cite{letter} introduces a collaborative contrastive loss and a diversity loss during tokenization to push co-interacted items toward proximate codes.
%
ColaRec~\cite{colarec} derives item identifiers from a pretrained collaborative filtering model and aligns them with content representations through contrastive learning and an item-indexing task.
%
UNGER~\cite{unger} integrates semantic and collaborative knowledge into a unified code via joint optimization of cross-modality alignment and next-item prediction.
%
MSCGRec~\cite{mscgrec} treats collaborative features extracted from a sequential recommender as an additional modality alongside text and image, and proposes constrained sequence learning to restrict the output space during training.
%
CEMG~\cite{cemg} dynamically integrates visual and textual features under the guidance of collaborative signals before unified RQ-VAE tokenization; however, it fuses all modalities into a single representation before quantization, sacrificing per-modality discriminability.
%
These methods each inject collaborative signals into the code learning process; however, they either operate on a single content modality~\cite{letter,colarec}, compress all signals into one code at the cost of per-modality discriminability~\cite{unger,cemg}, or route collaborative signals through a separate quantization path that does not reshape the text or image code structure~\cite{mscgrec}.

Beyond code-space alignment, ETEGRec~\cite{etegrec} bridges the tokenizer-recommender gap through alternating optimization with distribution-level and preference-level alignment losses, enabling co-adaptation but still relying on single-modality content features with inference restricted to discrete code space.

CAGRec builds upon MACRec's cross-modal RQ-VAE and further proposes collaborative preference anchoring to align user representations with the collaborative-semantic space.
%
It anchors multi-modal code learning with collaborative embeddings, structurally reshaping both text and image code spaces toward behavioral proximity while preserving per-modality paths, and leverages the resulting collaborative-semantic space for continuous reranking at inference.

%%%%%%%%%%%%%%%%%%%%%%%%%%%%%%%%%%%%%%%%%%%%%%%%%%%%%%%%%%%%%%%%%%%%%%%%
%% METHODOLOGY
%%%%%%%%%%%%%%%%%%%%%%%%%%%%%%%%%%%%%%%%%%%%%%%%%%%%%%%%%%%%%%%%%%%%%%%%

% ---- CAGRec 总框架图 ----
\begin{figure*}[t]
    \centering
    \includegraphics[width=\linewidth]{Model_CAGRec.pdf}
    \caption{Overall architecture of CAGRec. \textbf{(a)}~Collaborative-anchored cross-modal quantization: collaborative embeddings are concatenated with text and image features as anchored inputs for dual-path RQ-VAE, with $\mathcal{L}_{\mathrm{CAR}}$ reinforcing collaborative proximity in the latent space. \textbf{(b)}~Generative recommendation with CPA: a T5 model predicts target codes ($\mathcal{L}_{\mathrm{code}}$), and $\mathcal{L}_{\mathrm{CPA}}$ aligns the decoder's user representation to the target's collaborative-semantic embedding from the frozen RQ-VAE. \textbf{(c)}~Semantic-aware reranking: beam candidates are reranked by combining generation likelihood with collaborative-semantic cosine similarity.}
    \Description{Architecture diagram of CAGRec: (a) collaborative-anchored cross-modal RQ-VAE with collaborative-aware regularization; (b) T5-based recommendation training with collaborative preference anchoring; (c) semantic-aware reranking combining beam scores with cosine similarity.}
    \label{fig:model}
\end{figure*}

\section{Methodology}

\noindent\textbf{Problem formulation.}
Let $\mathcal{U}$ and $\mathcal{I}$ denote the user and item sets.
%
Each item $i \in \mathcal{I}$ is associated with a text description and an image.
%
For user $u$ with interaction history $[i_1, \ldots, i_n]$, the goal is to predict the next item $i_{n+1}$.

\noindent\textbf{Overview.}
Figure~\ref{fig:model} illustrates CAGRec, which injects collaborative signals at three stages.
%
To bridge the \emph{content-collaborative gap}, we anchor the cross-modal RQ-VAE with collaborative embeddings from user interaction sequences, steering code spaces toward behavioral structure~(\S\ref{sec:quant}).
%
The anchored RQ-VAE decoder gives rise to a \emph{collaborative-semantic space} that jointly encodes behavioral and content signals (formally defined in \S\ref{sec:quant}).
%
During training, we further propose Collaborative Preference Anchoring~(CPA) to ground user representations in this space~(\S\ref{sec:cpa}).
%
To address the \emph{discrete-continuous gap} at inference, we design a semantic-aware reranking strategy that re-evaluates candidates in the collaborative-semantic space~(\S\ref{sec:rerank}).

% ==================================================================
% 3.1 Collaborative-Anchored Cross-Modal Item Quantization
% ==================================================================
\subsection{Collaborative-Anchored Cross-Modal Item Quantization}
\label{sec:quant}

We build upon the cross-modal RQ-VAE framework of MACRec~\cite{macrec}, which introduces cross-modal contrastive learning during residual quantization.
%
However, its code spaces are built purely from content features: items with similar interaction patterns but distinct text or images end up with distant codes, forcing the downstream generative model to predict across a misaligned code structure.
%
Our key extension is to anchor the quantization input with collaborative embeddings derived from user behavior, so that the code structure respects behavioral proximity from the outset.

% ---- 协同嵌入与锚定输入 ----
\subsubsection{Collaborative Embedding and Anchored Input}

We pretrain a SASRec~\cite{sasrec} model on user interaction sequences and use its learned item embedding table to obtain a collaborative embedding $\mathbf{c}_i \in \mathbb{R}^{d_c}$ for each item~$i$.
% 
This embedding encodes item co-occurrence patterns learned from sequential context: items that frequently appear within similar interaction sequences receive proximate representations, regardless of their content features.
% 
We also encode the text description and image of item~$i$ into $\mathbf{t}_i \in \mathbb{R}^{d_t}$ and $\mathbf{v}_i \in \mathbb{R}^{d_v}$ using a pretrained sentence encoder and a vision encoder~\cite{clip}, respectively.

We construct the \emph{collaborative-anchored} input for each modality by concatenation:
\begin{equation}
    \mathbf{s}_i^t = [\mathbf{c}_i \,;\, \mathbf{t}_i] \in \mathbb{R}^{d_c+d_t}, \quad
    \mathbf{s}_i^v = [\mathbf{c}_i \,;\, \mathbf{v}_i] \in \mathbb{R}^{d_c+d_v},
\end{equation}
where $[\cdot\,;\,\cdot]$ denotes concatenation. For notational brevity, we omit the item subscript~$i$ when clear from context.
%
The shared $\mathbf{c}_i$ acts as a cross-modal anchor: both modality paths receive the same collaborative embedding, so items with similar co-occurrence patterns share an identical collaborative component in the input of \emph{both} paths, implicitly aligning the two encoders toward a common behavioral structure while the modality-specific component ($\mathbf{t}_i$ or $\mathbf{v}_i$) retains per-modality discriminability.

% ---- 双路 RQ-VAE ----
\subsubsection{Dual-Path Cross-Modal RQ-VAE}

We adopt a dual-path RQ-VAE architecture.
% 
Each path has an MLP encoder, an $L$-layer residual quantizer, and an MLP decoder.
% 
The encoders map the anchored inputs into a shared latent dimension:
%
\begin{equation}
    \mathbf{z}^t = \mathrm{Enc}_t(\mathbf{s}^t), \quad
    \mathbf{z}^v = \mathrm{Enc}_v(\mathbf{s}^v),
\end{equation}

where $\mathbf{z}^t, \mathbf{z}^v \in \mathbb{R}^{d_e}$ serve as the initial residuals $\mathbf{r}_1^t = \mathbf{z}^t$, $\mathbf{r}_1^v = \mathbf{z}^v$.
%
At the $l$-th RQ layer, each modality maintains an independent codebook $\mathcal{E}_l^{m} = \{\mathbf{e}_{l,k}^{m}\}_{k=1}^{M}$ ($m \in \{t,v\}$), where $M$ is the codebook size.
% 
Quantization selects the nearest codebook entry and updates the residual:
\begin{equation}
    c_l^{m} = \arg\min_k \|\mathbf{r}_l^m - \mathbf{e}_{l,k}^m\|_2, \quad
    \mathbf{r}_{l+1}^m = \mathbf{r}_l^m - \mathbf{e}_{l,c_l^m}^m.
\end{equation}
The code tuple $(c_1^m, \ldots, c_L^m)$ becomes the semantic ID for item~$i$ under modality~$m$.

Following MACRec~\cite{macrec}, we apply two cross-modal objectives during quantization to exploit inter-modal complementarity.
%
Starting from layer~$L_c$, a cross-modal contrastive loss applies InfoNCE~\cite{oord2018representation} on the RQ residuals using the other modality's K-means pseudo-labels as positive-pair indicators:
\begin{equation}
    \mathcal{L}_{\mathrm{con}}^l = \mathcal{L}_{\mathrm{con}}^{l,\, v \to t} + \mathcal{L}_{\mathrm{con}}^{l,\, t \to v}, \quad l \geq L_c,
\end{equation}
encouraging codebooks to capture complementary structures and reducing codebook collapse.
%
A reconstruction alignment loss $\mathcal{L}_{\text{align}}$ applies bidirectional InfoNCE on the quantized latents:
\begin{equation}
    \hat{\mathbf{z}}^m = \sum_{l=1}^{L}\mathbf{e}_{l,c_l^m}^m,
\end{equation}
to keep the text-path and image-path representations of the same item close.
%
Each path has a separate decoder that reconstructs the collaborative-anchored input: $\hat{\mathbf{s}}^m = \mathrm{Dec}_m(\hat{\mathbf{z}}^m) \in \mathbb{R}^{d_c+d_m}$.
%
We refer to this output space as the \emph{collaborative-semantic space}; it jointly encodes collaborative and content information and serves as the continuous target for both CPA~(\S\ref{sec:cpa_loss}) and reranking~(\S\ref{sec:rerank}).
%
We refer readers to~\citet{macrec} for the full formulation of these two losses.

% ---- 协同感知正则化 ----
\subsubsection{Collaborative-Aware Regularization}

Concatenation anchors the \emph{input} with collaborative structure; we further regularize the \emph{latent space} so that the learned representations respect collaborative proximity before quantization.
%
Inspired by the collaborative regularization in LETTER~\cite{letter}, we apply a contrastive loss on the encoder outputs $\mathbf{z}^t,\mathbf{z}^v$ \emph{before} quantization.
%
Unlike LETTER, which aligns post-quantization embeddings with CF embeddings in a single-modality setting via cross-space contrastive learning, our regularization operates in the pre-quantization latent space of both modality paths using collaborative neighbor relationships.
%
For each item~$i$, we precompute a set of collaborative neighbors $\mathcal{N}(i)$ as items with the highest cosine similarity in the SASRec embedding space.
%
For each modality $m \in \{t,v\}$, we apply InfoNCE where items sharing collaborative neighbors form positive pairs:
\begin{equation}
    \mathcal{L}_{\text{CAR}}^m = -\frac{1}{|\mathcal{P}^m|}\sum_{(i,j)\in\mathcal{P}^m}\log\frac{\exp(\cos(\mathbf{z}_i^m,\mathbf{z}_j^m)/\tau_c)}{\sum_{k}\exp(\cos(\mathbf{z}_i^m,\mathbf{z}_k^m)/\tau_c)},
\end{equation}
where $\mathcal{P}^m$ is the set of positive pairs $(i,j)$ with $j \in \mathcal{N}(i)$ within the current batch, and $\tau_c$ is the temperature.
%
This loss synergizes with collaborative anchoring: the concatenated collaborative embedding provides the initial behavioral structure in the input, while $\mathcal{L}_{\text{CAR}}$ reinforces this structure in the learned latent space, encouraging collaboratively related items to cluster before code assignment.

% ---- 量化阶段训练目标 ----
\subsubsection{Item Quantization Objective}

The overall pretraining objective for the dual-path RQ-VAE is:
\begin{equation}
    \mathcal{L}_{\text{quant}} = \mathcal{L}_{\text{recon}} + \lambda_q\,\mathcal{L}_{\text{vq}} + \sum_{l=L_c}^{L-1}\lambda_{\text{con}}^l\,\mathcal{L}_{\text{con}}^l + \lambda_a\,\mathcal{L}_{\text{align}} + \lambda_r\,\mathcal{L}_{\text{CAR}},
    \label{eq:quant}
\end{equation}
where $\mathcal{L}_{\text{recon}} = \|\mathbf{s}^t - \hat{\mathbf{s}}^t\|_2^2 + \|\mathbf{s}^v - \hat{\mathbf{s}}^v\|_2^2$ is the reconstruction loss, $\mathcal{L}_{\text{vq}}$ is the standard VQ commitment loss~\cite{image_rqvae}, $\mathcal{L}_{\text{CAR}} = \mathcal{L}_{\text{CAR}}^t + \mathcal{L}_{\text{CAR}}^v$ is the collaborative-aware regularization summed over both modalities, and $\lambda_q$, $\lambda_{\text{con}}^l$, $\lambda_a$, $\lambda_r$ are trade-off hyperparameters.
% 
When multiple items are mapped to the same code tuple (ID conflict), we follow~\citet{mql} and reassign conflicting items to their next-nearest codebook entries, ensuring a one-to-one mapping between items and codes.

% ==================================================================
% 3.2 Generative Recommendation with CPA
% ==================================================================
\subsection{Generative Recommendation with Collaborative Preference Anchoring}
\label{sec:cpa}

After pretraining the RQ-VAE, each item has a text semantic ID and an image semantic ID.
% 
We train a T5-based~\cite{t5} encoder-decoder model that autoregressively generates the next item's code given the user's interaction history.
%
Beyond the standard code prediction loss, we propose \emph{Collaborative Preference Anchoring}~(CPA) to ground the recommender's user representations in the continuous collaborative-semantic space produced by the anchored RQ-VAE decoder~(\S\ref{sec:quant}).
%
Because this space encodes both behavioral and content signals by construction, CPA implicitly encourages the recommender to capture collaborative proximity that pure content alignment would miss.

% ---- 双路序列推荐 ----
\subsubsection{Dual-Path Code Generation}

The text and image paths share the same T5 backbone but maintain path-specific token embedding tables and adapter layers.
% 
Given a user's interaction history encoded as a code sequence $\mathbf{x}^m$ for modality~$m$, the T5 encoder processes $\mathbf{x}^m$ and the decoder predicts the target code token by token:
\begin{equation}
    \mathcal{L}_{\text{code}}^m = -\sum_{l=1}^{L} \log P_\theta(c_l^m \mid c_{<l}^m,\, \mathbf{x}^m).
\end{equation}
%
Following MACRec~\cite{macrec}, we adopt six auxiliary training tasks (item-to-image, image-to-item, etc.) alongside the primary sequential recommendation tasks to enhance the model's cross-modal understanding.

% ---- CPA ----
\subsubsection{Collaborative Preference Anchoring (CPA)}
\label{sec:cpa_loss}

Code-level supervision treats all mismatched codes as equally wrong; CPA directly aligns the decoder's output with continuous collaborative-semantic embeddings, recovering fine-grained similarity that discretization discards.

For the target item, the RQ-VAE reconstructs its anchored semantic embedding: $\hat{\mathbf{s}}_{\text{tgt}}^m = \mathrm{Dec}_m(\hat{\mathbf{z}}_{\text{tgt}}^m)$.
% 
For the user, the T5 decoder hidden state at the \textsc{[bos]} position (index 0) captures the global user preference before any target token is generated.
%
We project it into the same semantic space via a learnable adapter $f_{\text{dec}}^m$:
\begin{equation}
    \mathbf{d}^m = f_{\text{dec}}^m\!\bigl(\mathrm{T5\text{-}Dec}(\mathbf{x}^m)_{\textsc{[bos]}}\bigr).
\end{equation}
CPA minimizes the cosine distance between the user prediction $\mathbf{d}^m$ and the target reconstruction $\hat{\mathbf{s}}_{\text{tgt}}^m$:
\begin{equation}
    \mathcal{L}_{\text{CPA}}^m = 1 - \cos(\mathbf{d}^m,\, \hat{\mathbf{s}}_{\text{tgt}}^m),
    \label{eq:cpa}
\end{equation}
where $\cos(\cdot,\cdot)$ denotes cosine similarity.
% 
Because the RQ-VAE was trained to reconstruct the collaborative-anchored input~$\mathbf{s}^m$ (i.e., $[\mathbf{c};\,\mathbf{t}]$ or $[\mathbf{c};\,\mathbf{v}]$), the reconstructed embedding $\hat{\mathbf{s}}_{\text{tgt}}^m$ encodes both collaborative patterns and content features.
%
Compared with contrastive objectives that require in-batch negatives, cosine distance loss provides a direct pairwise alignment signal that is more stable in our setting where the target space is already well-structured by the collaborative-anchored RQ-VAE.

% ---- 训练目标 ----
\subsubsection{Training Objective}
\label{sec:train_obj}

The total recommendation loss for each modality combines code prediction and CPA:
\begin{equation}
    \mathcal{L}_{\text{rec}}^m = \mathcal{L}_{\text{code}}^m + \lambda_p\,\mathcal{L}_{\text{CPA}}^m,
    \label{eq:rec}
\end{equation}
where $\lambda_p$ is the CPA loss weight.
% 
The final loss sums over both modalities: $\mathcal{L}_{\text{rec}} = \mathcal{L}_{\text{rec}}^t + \mathcal{L}_{\text{rec}}^v$.

% ==================================================================
% 3.3 Semantic-aware Reranking
% ==================================================================
\subsection{Semantic-aware Reranking}
\label{sec:rerank}

Beam search ranks candidates by token-level log-probabilities, operating entirely in the discrete code space.
% 
This makes it blind to continuous semantic relationships---two items with nearly identical collaborative-semantic embeddings can receive very different beam scores if their codes diverge at early positions.
% 
We bridge this \emph{discrete-continuous gap} by re-evaluating candidates in the continuous collaborative-semantic space at inference.

% ---- 用户语义表征 ----
\subsubsection{User Representation in Semantic Space}

At inference, the decoder has not yet started generation, so we derive the user representation from the encoder output.
%
For each modality path, we mean-pool the T5 encoder output and project it with the CPA-trained adapter $f_{\text{dec}}^m$:
\begin{equation}
    \mathbf{u}^m = f_{\text{dec}}^m\!\bigl(\mathrm{MeanPool}(\mathrm{T5\text{-}Enc}(\mathbf{x}^m))\bigr) \in \mathbb{R}^{d_c+d_m}.
\end{equation}
Although $f_{\text{dec}}^m$ is trained on decoder hidden states, the T5 decoder's initial state is derived from the encoder output via cross-attention, so the two representations lie in a closely related subspace; the ablation results in Section~4 empirically support this approximation.

% ---- 重排分数 ----
\subsubsection{Reranking Score}

For each beam candidate with code tuple $(c_1^m, \ldots, c_L^m)$, we look up the corresponding item~$j$ via a code-to-item index (which is one-to-one after reassignment) and retrieve its collaborative-anchored semantic embedding $\mathbf{s}_j^m$.
%
Candidates whose codes do not map to any item are discarded.
% 
The final score combines generation likelihood with collaborative-semantic similarity:
\begin{equation}
    \mathrm{score}(j) = \underbrace{s_{\text{beam}}(j)}_{\text{generation likelihood}} \;+\; \beta\cdot\underbrace{\cos(\mathbf{u}^m,\, \mathbf{s}_j^m)}_{\text{collaborative-semantic similarity}},
    \label{eq:rerank}
\end{equation}
where $\beta$ controls the reranking strength.
%
Because beam search log-probabilities and cosine similarities live on different scales, $\beta$ is intentionally kept small so that semantic similarity acts as a tiebreaker among candidates with comparable generation likelihood rather than overriding the autoregressive ranking.

% ---- 多模态集成 ----
\subsubsection{Multimodal Ensemble}

Following~\citet{macrec}, we fuse text-path and image-path results by score aggregation after reranking.
% 
For items appearing in both paths, scores are averaged; items appearing in only one path contribute directly.
% 
The fused ranking list is used for final evaluation.

%% NOTE: table* placed early so LaTeX can position it on an earlier page
\begin{table*}[t!]
\caption{Performance comparison on three datasets. Bold and underline indicate the best and second-best results, respectively. $\dagger$: reproduced by us under identical settings; all other baseline results are cited from~\cite{macrec}.}
\label{tab:main}
\centering
\resizebox{\textwidth}{!}{%
\begin{tabular}{l|l|*{10}{C}}
\toprule
\multicolumn{2}{c|}{} & \multicolumn{3}{c}{\textit{Traditional Seq.}} & \textit{MM Seq.} & \multicolumn{2}{c}{\textit{Single-modal Gen.}} & \multicolumn{4}{c}{\textit{MM Gen.}} \\
\cmidrule(lr){3-5} \cmidrule(lr){6-6} \cmidrule(lr){7-8} \cmidrule(lr){9-12}
Dataset & Metrics & BERT4Rec & SASRec & S$^3$-Rec & MISSRec & P5-CID & TIGER & VIP5 & MQL4GRec & MACRec$^\dagger$ & CAGRec\\
\midrule
\multirow{5}{*}{Instruments} 
& HR@1      & 0.0450 & 0.0318 & 0.0339 & 0.0723 & 0.0512 & 0.0754 & 0.0737 & 0.0763 & \underline{0.0804} & \textbf{0.0830} \\
& HR@5      & 0.0856 & 0.0946 & 0.0937 & 0.1089 & 0.0839 & 0.1007 & 0.0892 & 0.1058 & \underline{0.1112} & \textbf{0.1161} \\
& HR@10     & 0.1081 & 0.1233 & 0.1123 & \underline{0.1361} & 0.1119 & 0.1221 & 0.1071 & 0.1291 & 0.1349 & \textbf{0.1401} \\
& NDCG@5    & 0.0667 & 0.0654 & 0.0693 & 0.0797 & 0.0678 & 0.0882 & 0.0815 & 0.0902 & \underline{0.0962} & \textbf{0.0995} \\
& NDCG@10   & 0.0739 & 0.0746 & 0.0743 & 0.0880 & 0.0704 & 0.0950 & 0.0872 & 0.0997 & \underline{0.1038} & \textbf{0.1072} \\
\midrule
\multirow{5}{*}{Arts}
& HR@1      & 0.0289 & 0.0212 & 0.0172 & 0.0479 & 0.0421 & 0.0532 & 0.0474 & \textbf{0.0626} & \underline{0.0611} & \textbf{0.0626} \\
& HR@5      & 0.0697 & 0.0951 & 0.0739 & \underline{0.1021} & 0.0713 & 0.0894 & 0.0704 & 0.0997 & 0.0990 & \textbf{0.1027} \\
& HR@10     & 0.0922 & 0.1250 & 0.1030 & 0.1321 & 0.0994 & 0.1167 & 0.0959 & \underline{0.1254} & 0.1245 & \textbf{0.1296} \\
& NDCG@5    & 0.0502 & 0.0610 & 0.0511 & 0.0699 & 0.0607 & 0.0718 & 0.0586 & \underline{0.0816} & 0.0802 & \textbf{0.0830} \\
& NDCG@10   & 0.0575 & 0.0706 & 0.0630 & 0.0815 & 0.0662 & 0.0806 & 0.0635 & \underline{0.0898} & 0.0884 & \textbf{0.0916} \\
\midrule
\multirow{5}{*}{Games}
& HR@1      & 0.0115 & 0.0069 & 0.0136 & \underline{0.0201} & 0.0169 & 0.0166 & 0.0173 & 0.0200 & 0.0192 & \textbf{0.0202} \\
& HR@5      & 0.0426 & 0.0587 & 0.0527 & \underline{0.0674} & 0.0532 & 0.0523 & 0.0480 & 0.0645 & 0.0647 & \textbf{0.0675} \\
& HR@10     & 0.0725 & 0.0985 & 0.0903 & \underline{0.1048} & 0.0824 & 0.0857 & 0.0758 & 0.1007 & 0.1041 & \textbf{0.1078} \\
& NDCG@5    & 0.0270 & 0.0333 & 0.0351 & 0.0385 & 0.0331 & 0.0345 & 0.0328 & \underline{0.0421} & 0.0418 & \textbf{0.0437} \\
& NDCG@10   & 0.0366 & 0.0461 & 0.0468 & 0.0499 & 0.0454 & 0.0453 & 0.0418 & 0.0538 & \underline{0.0545} & \textbf{0.0565} \\
\bottomrule
\end{tabular}}%
\end{table*}

%%%%%%%%%%%%%%%%%%%%%%%%%%%%%%%%%%%%%%%%%%%%%%%%%%%%%%%%%%%%%%%%%%%%%%%%
%% EXPERIMENTS
%%%%%%%%%%%%%%%%%%%%%%%%%%%%%%%%%%%%%%%%%%%%%%%%%%%%%%%%%%%%%%%%%%%%%%%%

\section{Experiment}

We evaluate CAGRec along four research questions:
\textbf{RQ1}:~How does CAGRec compare to state-of-the-art baselines?
\textbf{RQ2}:~How does each component contribute?
\textbf{RQ3}:~How does the fusion strategy for collaborative embeddings affect performance?
\textbf{RQ4}:~How sensitive is CAGRec to key hyperparameters?

\subsection{Experimental Setup}

\noindent$\bullet$~\textbf{Datasets.}
We use three categories from the Amazon Product Reviews (2018) dataset~\cite{ni2019justifying}: ``Musical Instruments'' (Instruments), ``Arts, Crafts and Sewing'' (Arts), and ``Video Games'' (Games).
% 
Each category contains user reviews with ratings, product metadata (title, description, brand), and product images.
% 
Following prior work~\cite{tiger,macrec}, we apply 5-core filtering (retaining users and items with at least 5 interactions) and chronologically sort each user's interactions.
% 
Table~\ref{dataset} summarizes the statistics after preprocessing.

\begin{table}[h]
\caption{Statistics of the datasets. \textbf{Avg. len} represents the average length of item sequences.}
\label{dataset}
\small
\centering
\setlength\tabcolsep{2pt}
\begin{tabular}{lccccc}
\toprule
\textbf{Datasets} & \textbf{\#Users} & \textbf{\#Items} & \textbf{\#Interactions} & \textbf{Sparsity} & \textbf{Avg. len} \\
\midrule
Instruments & 17112 & 6250  & 136226  & 99.87\% & 7.96 \\
Arts       & 22171 & 9416  & 174079  & 99.92\% & 7.85 \\
Games      & 42259 & 13839 & 373514  & 99.94\% & 8.84 \\
\bottomrule
\end{tabular}
\end{table}

\noindent$\bullet$~\textbf{Baselines.}
We compare with nine representative methods spanning four categories.
%
\textit{(i) Traditional sequential recommenders}: \textbf{SASRec}~\cite{sasrec} uses unidirectional self-attention over item ID sequences; \textbf{BERT4Rec}~\cite{sun2019bert4rec} extends this with bidirectional masked prediction; \textbf{S$^3$-Rec}~\cite{zhou2020s3} further introduces self-supervised pre-training with mutual information maximization on item attributes.
%
\textit{(ii) Multimodal sequential recommender}: \textbf{MISSRec}~\cite{wang2023missrec} pre-trains a multimodal item encoder with interest-aware contrastive learning for transferable recommendation.
%
\textit{(iii) Single-modal generative recommenders}: \textbf{P5-CID}~\cite{p5,hua2023index} formulates recommendation as a unified text-to-text task with collaborative indexing that assigns each item a meaningful ID; \textbf{TIGER}~\cite{tiger} quantizes text embeddings via RQ-VAE into hierarchical semantic IDs and generates them autoregressively.
%
\textit{(iv) Multimodal generative recommenders}: \textbf{VIP5}~\cite{vip5} extends the P5 framework with visual features through a multimodal adapter; \textbf{MQL4GRec}~\cite{mql} encodes text and image into a shared quantitative language with multi-level codebooks; \textbf{MACRec}~\cite{macrec} applies cross-modal contrastive learning during residual quantization to reduce codebook collisions across text and image paths.

\noindent$\bullet$~\textbf{Metrics.}
We report top-$k$ hit rate (HR@$K$) and normalized discounted cumulative gain (NDCG@$K$) for $K \in \{1, 5, 10\}$.
% 
Following~\citet{p5,hua2023index}, we adopt a leave-one-out protocol with full ranking over the entire item set.

\noindent$\bullet$~\textbf{Implementation Details.}
We obtain collaborative embeddings ($d_c\!=\!256$) from a SASRec~\cite{sasrec} model pretrained strictly on the training split (excluding validation and test interactions under the leave-one-out protocol), text features ($d_t\!=\!4096$) from LLaMA-7B~\cite{touvron2023llama} mean pooling, and image features ($d_v\!=\!768$) from ViT-L/14~\cite{clip}.
% 
For CAR, each item's collaborative neighbor set $\mathcal{N}(i)$ contains the top-$10$ items by cosine similarity in the SASRec embedding space.
%
For RQ-VAE, we set codebook size $M\!=\!256$ with $L\!=\!4$ levels, CAR weight $\lambda_r\!=\!2.0$, temperature $\tau_c\!=\!0.1$, cross-modal contrastive weights $\lambda_{\text{con}}^{0,1}\!=\!0$, $\lambda_{\text{con}}^{2,3}\!=\!0.1$, and reconstruction alignment weight $\lambda_a\!=\!0.001$.
% 
Following~\citet{tiger,mql}, we use T5~\cite{t5} as the backbone with 4 transformer layers, 6 attention heads, and dimension~64 per head.
% 
We train with six auxiliary tasks following MACRec~\cite{macrec}.
% 
CPA uses cosine distance with weight $\lambda_p\!=\!0.005$ and a linear adapter projecting from T5 hidden dimension to the $(d_c\!+\!d_m)$-dimensional collaborative-semantic space.
% 
We use AdamW (batch size 1024, learning rate $10^{-3}$), beam search with $B\!=\!20$, dual-path score ensemble, and semantic-aware reranking strength $\beta\!=\!0.1$.
%
Under the leave-one-out protocol, we use the second-to-last interaction of each user as the validation set for hyperparameter selection; all hyperparameters ($\lambda_r$, $\lambda_p$, $\lambda_a$, $\beta$) are tuned on the validation set and fixed for test evaluation.
%
All reported results are averaged over 5 independent runs with different random seeds.
% 
For MQL4GRec, we disable pretraining on additional-category datasets for a fair comparison.

% ====================================================================
\subsection{Overall Performance (RQ1)}
\label{sec:rq1}

Table~\ref{tab:main} compares all methods across three datasets.
%
We reproduce MACRec (marked with~$\dagger$) under its published settings; all other baseline results are cited from~\cite{macrec}.

\noindent$\bullet$~\textbf{CAGRec achieves the best performance on every metric and dataset.}
%
On NDCG@10, CAGRec surpasses the strongest baseline by $+2.0\%$--$3.7\%$ across three datasets with sparsity levels ranging from 99.87\% to 99.94\%.

\noindent$\bullet$~\textbf{Improvements are consistent across both recall and ranking metrics.}
%
Beyond the NDCG@10 gains, HR@10 improves by $+2.9\%$ to $+3.4\%$ across all three datasets, confirming that collaborative anchoring helps the autoregressive model surface a broader set of relevant candidates.

\noindent$\bullet$~\textbf{Among generative baselines, MQL4GRec and MISSRec are the strongest competitors.}
%
CAGRec matches MQL4GRec on Arts HR@1 (both 0.0626) while outperforming it on all other metrics and surpassing both competitors on NDCG across all datasets, indicating that collaborative anchoring provides benefits beyond richer content encodings alone.
%
On Instruments, CAGRec achieves HR@1$\,=\,$0.0830 ($+3.2\%$ over MACRec$^\dagger$); on Games the gain reaches $+5.2\%$.
%
The larger gain on Games---the sparsest dataset---suggests that collaborative anchoring is especially beneficial when content signals alone cannot discriminate among a large item set.

% ====================================================================
\subsection{Component Ablation Study (RQ2)}
\label{sec:rq2}

\begin{table}[t]
\caption{Component ablation on Instruments (6{,}250 items) and Games (13{,}839 items). Starting from the full CAGRec, we remove one component at a time.}
\label{tab:ablation}
\centering
\setlength\tabcolsep{4pt}
\small
\begin{tabular}{l|cc|cc}
\toprule
\multirow{2}{*}{Model} & \multicolumn{2}{c|}{Instruments} & \multicolumn{2}{c}{Games} \\
\cmidrule(lr){2-3}\cmidrule(lr){4-5}
 & HR@10 & NDCG@10 & HR@10 & NDCG@10 \\
\midrule
\textbf{CAGRec} (full)  & \textbf{0.1401} & \textbf{0.1072} & 0.1078 & 0.0565              \\
\;\;w/o Input Anchoring  & 0.1334          & 0.1035          & 0.1024       & 0.0545              \\
\;\;w/o CAR              & 0.1308          & 0.1012          & 0.0982       & 0.0532              \\
\;\;w/o CPA              & 0.1369          & 0.1057          & \textbf{0.1091}  & \textbf{0.0568}              \\
\;\;w/o Reranking        & 0.1379          & 0.1060          & 0.1057       & 0.0558              \\
\bottomrule
\end{tabular}
\end{table}

We ablate each component from the full CAGRec (Table~\ref{tab:ablation}).
%
We test on Instruments (6{,}250 items, 99.87\% sparsity) and Games (13{,}839 items, 99.94\% sparsity), which span the two extremes of item-set size and interaction sparsity in our benchmark (Table~\ref{dataset}); Arts lies between them on both dimensions, so the two endpoints suffice to reveal scale-dependent trends.

\noindent$\bullet$~\textbf{w/o Input Anchoring} removes collaborative embedding concatenation at the RQ-VAE input.
%
On Instruments, NDCG@10 drops by $-3.5\%$ and HR@10 by $-4.8\%$; on Games the degradation is comparable ($-3.5\%$ NDCG@10, $-5.0\%$ HR@10), confirming that feature-level anchoring is the primary channel for injecting collaborative knowledge and that the benefit persists as the item set scales.

\noindent$\bullet$~\textbf{w/o CAR} removes the contrastive constraint that clusters collaboratively related items before code assignment.
%
This produces the largest single-component degradation on both datasets: NDCG@10 falls by $-5.6\%$ on Instruments and $-5.8\%$ on Games; HR@10 drops by $-6.6\%$ and $-8.9\%$, respectively.
%
Without contrastive pressure the quantizer minimizes reconstruction loss freely, but the resulting codes lose collaborative-semantic structure, yielding the worst downstream performance among all ablation variants.
%
The larger degradation on Games indicates that a bigger item space amplifies the need for explicit latent-level collaborative regularization.

\noindent$\bullet$~\textbf{w/o CPA} removes the preference anchoring loss from T5 training.
%
On Instruments, NDCG@10 drops by $-1.4\%$ and HR@10 by $-2.3\%$, confirming that decoder-level alignment provides complementary refinement beyond code-level supervision.
%
On Games, removing CPA marginally improves both metrics ($+0.5\%$ NDCG@10), likely because the larger item set (13{,}839) and higher sparsity introduce more quantization error into the frozen RQ-VAE targets, offsetting CPA's alignment benefit.
%
This suggests CPA is most effective when the code space is compact enough for reconstruction targets to remain faithful.

\noindent$\bullet$~\textbf{w/o Reranking} sets $\beta\!=\!0$, ranking candidates purely by beam search log-probability.
%
On Instruments, NDCG@10 drops by $-1.1\%$ and HR@10 by $-1.6\%$; on Games, the drops are $-1.2\%$ and $-1.9\%$, respectively.
%
The consistent positive contribution across both datasets---in contrast to CPA's dataset-dependent behavior---indicates that continuous collaborative-semantic similarity is a robust complementary signal for refining discrete beam search outputs.
%
The larger relative gain on HR@10 than on NDCG@10 suggests that semantic reranking primarily rescues relevant items that beam search underranks at lower positions rather than improving the top-1 prediction.

\noindent Across both datasets, code-space construction dominates: CAR alone accounts for up to $-5.8\%$ NDCG@10 degradation (Games), and Input Anchoring follows closely ($-3.5\%$ on both).
%
Reranking provides consistent gains ($\sim\!1\%$ NDCG@10), while CPA's contribution is dataset-dependent, highlighting that the RQ-VAE stage is universally the most impactful.
%
Mapping back to the two challenges in Section~1: Input Anchoring and CAR jointly address the \emph{content-collaborative gap} by restructuring code spaces toward behavioral proximity, accounting for the majority of the performance gain; semantic-aware reranking addresses the \emph{discrete-continuous gap} by re-evaluating candidates in continuous space, providing a complementary but consistent improvement.

%% NOTE: figure* placed early so LaTeX positions it on an earlier page
\begin{figure*}[t]
    \centering
    \includegraphics[width=\linewidth]{fig_rq4_sensitivity.pdf}
    \caption{Hyperparameter sensitivity on Instruments. Each panel varies one hyperparameter while fixing the others at default values. Stars mark the optimal setting.}
    \Description{Four-panel line chart showing hyperparameter sensitivity of CAGRec on Instruments.}
    \label{fig:hyperparam}
\end{figure*}

% ====================================================================
\subsection{Collaborative Embedding Fusion Analysis (RQ3)}
\label{sec:rq3}

We compare four strategies for fusing collaborative embeddings with content features in the RQ-VAE.

\begin{table}[t]
\caption{Comparison of collaborative embedding fusion strategies on Instruments. \textbf{Coll.} reports the text / image codebook collision rate (\%, lower is better).}
\label{tab:fusion}
\centering
\setlength\tabcolsep{3pt}
\small
\begin{tabular}{l|c|cc|cc}
\toprule
\multirow{2}{*}{Fusion Strategy} & Coll.(\%) & \multicolumn{2}{c|}{HR} & \multicolumn{2}{c}{NDCG} \\
\cmidrule(lr){3-4}\cmidrule(lr){5-6}
 & text / image & @5 & @10 & @5 & @10 \\
\midrule
Proj.\ Addition   & 23.9 / 38.2 & 0.1106 & 0.1364 & 0.0957 & 0.1040 \\
Gating            & 73.3 / 90.1 & 0.1138 & 0.1367 & 0.0982 & 0.1056 \\
Cross-Attention   & 48.3 / 70.9 & 0.1107 & 0.1342 & 0.0963 & 0.1038 \\
\midrule
Concat (ours)     & \textbf{5.0 / 2.0} & \textbf{0.1156} & \textbf{0.1400} & \textbf{0.0992} & \textbf{0.1070} \\
\bottomrule
\end{tabular}
\end{table}

Table~\ref{tab:fusion} compares four fusion strategies on Instruments, all using the same SASRec collaborative embedding ($d_c\!=\!256$) and identical T5 configuration (including CPA).
%
A key architectural distinction separates our approach from the alternatives: concatenation operates at the \emph{raw input level}---the collaborative embedding is concatenated with the content feature before the align encoder, forming an extended input ($d_c\!+\!d_t$ or $d_c\!+\!d_v$ dimensions) that the encoder processes in full.
%
In contrast, the three parameterized alternatives all operate in a \emph{compressed intermediate space}: the align encoder first maps the content feature to a 768d representation, the collaborative embedding is separately projected to 768d via a learned linear layer, and the two 768d vectors are then combined:
%
(i)~\emph{Proj.\ Addition} adds the projected collaborative vector to the content representation element-wise;
%
(ii)~\emph{Gating} learns a sigmoid gate $\mathbf{g} = \sigma(\mathbf{W}_g[\mathbf{c}_{\mathrm{proj}};\,\mathbf{x}_{\mathrm{align}}])$ that determines the per-dimension mixture: $\mathbf{g} \odot \mathbf{c}_{\mathrm{proj}} + (1-\mathbf{g}) \odot \mathbf{x}_{\mathrm{align}}$;
%
(iii)~\emph{Cross-Attention} stacks the collaborative and content vectors as two tokens, applies multi-head self-attention (4 heads), and takes the content-position output with a residual connection.
%
Because all three alternatives project both signals into a shared 768d space---compressing the content features from their original dimensionality (e.g., $4096$d for text) while up-projecting the $256$d collaborative embedding---the two signals are forced to share the same dimensions and the encoder must learn to disentangle them.
%
In contrast, concatenation preserves each signal in its own subspace of the full-dimensional input, giving the MLP encoder strictly more capacity for flexible combination.

\noindent$\bullet$~\textbf{Concatenation achieves the best performance with the lowest collision rate.}
%
Our approach outperforms all parameterized fusion strategies across every metric (Table~\ref{tab:fusion}), with an even larger gap in codebook collision rate: 5.0\%/2.0\% (text/image) versus 23.9\%--73.3\% / 38.2\%--90.1\% for the alternatives.

\noindent$\bullet$~\textbf{Parameterized fusion destabilizes codebook learning.}
%
Gating exhibits the most extreme collapse (73.3\%/90.1\% collision), indicating that the adaptive gate tends to suppress content signals and converge to a degenerate solution during RQ-VAE training.
%
Cross-Attention (48.3\%/70.9\%) and Proj.\ Addition (23.9\%/38.2\%) also suffer significant collision inflation.
%
All three share a common failure mode: learnable parameters that mix collaborative and content signals \emph{before} quantization create unstable gradients during codebook updates, causing many items to collapse to identical codes.

\noindent$\bullet$~\textbf{Collision rate and downstream performance are not perfectly correlated.}
%
Despite having the worst collision rate, Gating outperforms Proj.\ Addition and Cross-Attention on NDCG@10, because index reassignment (\S\ref{sec:quant}) partially compensates for codebook collapse.
%
Nevertheless, concatenation's inherently low collision rate avoids the need for such post-hoc correction, yielding codes that directly reflect the intended collaborative-content structure.

% ====================================================================
\subsection{Hyperparameter Sensitivity (RQ4)}
\label{sec:rq4}

Figure~\ref{fig:hyperparam} reports HR@10 and NDCG@10 on Instruments as we vary four key hyperparameters one at a time.

\noindent$\bullet$~\textbf{CAR weight $\lambda_r$ (Fig.~\ref{fig:hyperparam}a).}
%
Both metrics peak at $\lambda_r\!=\!2.0$, where NDCG@10 reaches $0.1072$.
%
When $\lambda_r$ is too small ($0.5$--$1.0$), the contrastive signal fails to shape the codebook toward collaborative structure, and NDCG@10 drops by $2.0$--$2.3\%$.
%
When $\lambda_r\!=\!5.0$, excessive contrastive pressure degrades reconstruction quality, reducing NDCG@10 by $1.3\%$.
%
Among all four hyperparameters, $\lambda_r$ exhibits the highest sensitivity, consistent with CAR being the most critical loss component (cf.\ RQ2).

\noindent$\bullet$~\textbf{CPA weight $\lambda_p$ (Fig.~\ref{fig:hyperparam}b).}
%
$\lambda_p\!=\!0.005$ yields the best NDCG@10 (0.1072).
%
At $\lambda_p\!=\!0.002$ the alignment signal is too weak; at $\lambda_p\!=\!0.01$ the decoder over-optimizes cosine alignment at the cost of code prediction accuracy.
%
The sensitivity profile resembles $\lambda_r$, with a clear inverted-U shape.

\noindent$\bullet$~\textbf{Alignment weight $\lambda_a$ (Fig.~\ref{fig:hyperparam}c).}
%
$\lambda_a$ controls the reconstruction--contrastive alignment term in the RQ-VAE.
%
NDCG@10 peaks at $\lambda_a\!=\!0.001$ (0.1072) and degrades mildly at both extremes ($-0.7\%$ at $10^{-4}$, $-1.3\%$ at $10^{-2}$).
%
HR@10 remains nearly flat across all three values (0.137--0.140), indicating that $\lambda_a$ primarily affects ranking precision rather than recall.

\noindent$\bullet$~\textbf{Reranking strength $\beta$ (Fig.~\ref{fig:hyperparam}d).}
%
Both metrics exhibit a clear inverted-U pattern.
%
Without reranking ($\beta\!=\!0$), HR@10 is 0.1379 and NDCG@10 is 0.1060.
%
At the optimal $\beta\!=\!0.1$, both metrics peak (HR@10$\,=\,$0.1401, $+1.6\%$; NDCG@10$\,=\,$0.1072, $+1.1\%$), confirming that collaborative-semantic similarity provides an effective complementary signal for refining generation scores.
%
When $\beta$ increases to $0.2$, performance drops (HR@10$\,=\,$0.1390, NDCG@10$\,=\,$0.1066), as the collaborative prior begins to override the well-calibrated generation likelihood.
%
A moderate reranking weight recovers items that are behaviorally relevant but ranked lower by the generator; excessive weight degrades precision.

% ====================================================================
% [RQ5 - Visualization Analysis: moved to appendix]
% \subsection{Visualization Analysis (RQ5)}
% \label{sec:rq5}
%
% \begin{figure}[t]
%     \centering
%     \subfloat[Text (by Brand)]{\includegraphics[width=0.48\linewidth]{fig1_a.pdf}\label{fig:tsne-a}}
%     \hfill
%     \subfloat[Collab (by Brand)]{\includegraphics[width=0.48\linewidth]{fig1_b.pdf}\label{fig:tsne-b}}
%     \\[2pt]
%     \subfloat[Image (by Category)]{\includegraphics[width=0.48\linewidth]{fig1_c.pdf}\label{fig:tsne-c}}
%     \hfill
%     \subfloat[Collab (by Category)]{\includegraphics[width=0.48\linewidth]{fig1_d.pdf}\label{fig:tsne-d}}
%     \caption{t-SNE visualization on Instruments. Text embeddings cluster items by brand~(a) and image embeddings by product type~(c). The same items, with identical coloring, scatter across the collaborative space~(b)(d). Content similarity does not equal collaborative similarity.}
%     \Description{Four t-SNE visualizations in a two-by-two grid comparing content and collaborative embedding spaces.}
%     \label{fig:tsne}
% \end{figure}
%
% Figure~\ref{fig:tsne} visualizes item embeddings on Instruments via t-SNE.
% %
% Text embeddings cluster items by brand~(a) and image embeddings by product type~(c), yet the same items scatter across the collaborative space~(b)(d).
% %
% Content proximity and behavioral proximity are two distinct structures---the core observation that motivates collaborative anchoring.

%%%%%%%%%%%%%%%%%%%%%%%%%%%%%%%%%%%%%%%%%%%%%%%%%%%%%%%%%%%%%%%%%%%%%%%%
%% CONCLUSION
%%%%%%%%%%%%%%%%%%%%%%%%%%%%%%%%%%%%%%%%%%%%%%%%%%%%%%%%%%%%%%%%%%%%%%%%

\section{Conclusion and Limitation}

We identify the mismatch between content-built code spaces and collaborative structure in multimodal generative recommendation, and propose CAGRec to inject collaborative signals at three complementary stages: input-level concatenation, latent-level regularization (CAR), and training-level preference anchoring (CPA).
%
Ablation studies confirm that each stage contributes complementary gains: input anchoring and CAR jointly reshape the code space, while CPA and semantic-aware reranking refine ranking precision.
% 
Experiments on three Amazon datasets show consistent improvements over all baselines, with NDCG@10 gains of $+2.0\%$ to $+3.7\%$ over the strongest competitor on each dataset.

Despite these results, collaborative anchoring relies on pretrained sequential embeddings, leaving cold-start items without effective anchoring.
%
Moreover, the two-stage pipeline freezes collaborative embeddings during recommendation training; end-to-end tokenizer-recommender co-training~\cite{etegrec} could enable joint adaptation and further improve code quality.

%%%%%%%%%%%%%%%%%%%%%%%%%%%%%%%%%%%%%%%%%%%%%%%%%%%%%%%%%%%%%%%%%%%%%%%%
%% REFERENCES
%%%%%%%%%%%%%%%%%%%%%%%%%%%%%%%%%%%%%%%%%%%%%%%%%%%%%%%%%%%%%%%%%%%%%%%%

\bibliographystyle{ACM-Reference-Format}
\bibliography{references}

\end{document}